% !PRINTED 703
\weberpage{703}
%<<CONT
the theory of infinite series a sum of terms of the
form
\[
\left(\frac{1}{N(\mathfrak{p})^k}\,\frac{1}{N(\mathfrak{p}')^{k'}}\,\frac{1}{N(\mathfrak{p}'')^{k''}}\cdots\right)^s = \left(\frac{1}{N(\mathfrak{p}^k\mathfrak{p}'^{k'}\mathfrak{p}''^{k''}\dots)}\right)^s,
\]
all of which are contained in the form $N(\mathfrak{a})^{-s}$, and each such
term arises once and only once. Accordingly we have
\weq{5}{\varPhi(s) = \prod \frac{1}{1-N(\mathfrak{p})^{-s}}.}

If therefore $\mathfrak{p}_1, \mathfrak{p}_2, \dots, \mathfrak{p}_e$ are the prime factors of $p$ that are
distinct from one another, and $f_1, f_2, \dots, f_e$ their degrees (\S.\,143), then
we obtain
\weq{6}{\varPhi(s) = \overset{p}{\prod} \frac{1}{(1-p^{-sf_1})(1-p^{-sf_2})\dots(1-p^{-sf_e})},}
in which the infinite product $\prod$ is to be extended over all natural
 primes $p$, and by (3), (4)
\weq{7}{gh = \mathop{\mathrm{Lim}}\limits_{s=1}\,(s-1)\varPhi(s).}

From this a conclusion may be drawn that is important for
manifold applications:

If from the development
\[
1 + \frac{1}{p^s} + \frac{1}{p^{2s}} + \cdots = \frac{1}{1-p^{-s}}
\]
we form the product for all primes $p$, then we obtain:
\weq{8}{\overset{p}{\prod}\frac{1}{1-p^{-s}} = \overset{n}{\sum}\frac{1}{n^s},}
in which the sum on the right side extends over all natural
numbers $n$; and this product therefore has a finite value for all
values of $s$ that are greater than 1. The reciprocal
value, namely the product
\weq{9}{P = \overset{p}{\prod}(1-p^{-s}),}
whose factors are all less than 1, has therefore, as
long as $s>1$, a value different from zero. This
property is now preserved if, in forming the
product $P$, we take account of only a part of all the primes $p$,
because by the omission of arbitrary factors the product is only
increased, without ever exceeding unity.
